\documentclass[aps,prc,onecolumn,superscriptaddress,nofootinbib,10pt]{revtex4-2}
\usepackage{amsmath,amssymb}
\usepackage{booktabs}
\usepackage[colorlinks=true,linkcolor=blue,citecolor=blue,urlcolor=blue]{hyperref}

\newcommand{\nB}{n_B}
\newcommand{\nC}{n_C}
\newcommand{\nS}{n_S}
\newcommand{\nq}{n_q}
\newcommand{\muB}{\mu_B}
\newcommand{\muC}{\mu_C}
\newcommand{\muS}{\mu_S}
\newcommand{\munue}{\mu_{\nu_e}}
\newcommand{\Lam}{\Lambda}
\newcommand{\GS}{G_S}
\newcommand{\GD}{G_D}
\newcommand{\GV}{G_V}
\newcommand{\trace}{\mathrm{tr}}

\begin{document}

\title{Three-flavour NJL with colour superconductivity \\
       --- as implemented in \texttt{eos/njl}}

\author{M.~Guerrini}
\affiliation{University of Ferrara}

\begin{abstract}
The \texttt{eos/njl} subpackage implements the three-flavour
Nambu--Jona-Lasinio model of dense quark matter with a diquark condensate: the
scalar channel that breaks chiral symmetry, the 't~Hooft determinant that ties
the flavours together, the colour-antitriplet diquark channel that condenses
into a colour superconductor, and a vector channel whose repulsion controls the
high-density stiffness. Almost nothing in it is an input: the constituent
masses come out of a gap equation, the effective bag constant is a derived
vacuum pressure difference, and \emph{which} colour-superconducting pattern the
matter is in is an outcome chosen by free energy rather than a declaration.
This note states the Lagrangian, the cut medium integrals and the surface term
that distinguishes the two pressure forms, the Dirac sea in closed form, the
gap equations for masses and for gaps, the Bogoliubov--de~Gennes construction
and the Hellmann--Feynman kernels it supplies, the density-dependent vector
coupling and the rearrangement self-energy it owes, the residual of every
equilibrium mode as the solver assembles it, and the phase-adapter surface the
hybrid engine consumes. It is written from the code, as its specification:
every equation here is asserted either by
\texttt{eos/njl/verify/run\_full\_check.py}, by \texttt{test/njl}, or --- for
the shared pairing machinery --- by \texttt{test/general/test\_pairing.py}.
\end{abstract}

\maketitle

\section{Conventions}
\label{sec:conventions}

Units are natural, $\hbar = c = k_B = 1$, with MeV throughout: momenta, masses
and chemical potentials in MeV, densities in MeV$^3$, and $\Omega$, $P$ and
$\varepsilon$ in MeV$^4$. Every public boundary of the package converts to
fm-based units, densities in fm$^{-3}$ and $P$, $\varepsilon$ in MeV\,fm$^{-3}$,
through $(\hbar c)^3 = (197.3269804\ \mathrm{MeV\,fm})^3$.

\paragraph{Modes.} There are nine colour--flavour modes $j = (f,a)$ with
$f \in \{u,d,s\}$ and $a \in \{r,g,b\}$, indexed \emph{flavour-major},
$j = 3 i_f + i_a$. The spin degeneracy of one mode is $g = 2$; the Dirac sea of
one flavour carries $g_{\rm sea} = 2N_c = 6$, since the vacuum is not resolved
by colour.

\paragraph{Charges.} $q_u = +\tfrac23$, $q_d = q_s = -\tfrac13$, and
\emph{strangeness is $s_s = +1$ per $s$ quark}, the opposite of the PDG sign
and used consistently throughout this repository. $C$ is the charge of
strongly-interacting matter only; the leptons are excluded from it and enter
through the separate condition of total electric neutrality.

\paragraph{Colour generators.} $(T_3)_{r,g,b} = (+\tfrac12,-\tfrac12,0)$ and
\begin{equation}
  (T_8)_{r,g,b} = \left(\tfrac13,\ \tfrac13,\ -\tfrac23\right)
  = \frac{\lambda_8}{\sqrt3} .
\end{equation}
\emph{Three} normalisations of $T_8$ are in circulation and mixing them
corrupts $\mu_8$ by factors of $1.15$ to $1.7$. R\"uster \textit{et
al.}~\cite{Ruester2005}, Pagliara and Schaffner-Bielich~\cite{Pagliara2008} and
Kunkel \textit{et al.}~\cite{Kunkel2026} use the halved Gell-Mann
$\sqrt3\,T_8 = \mathrm{diag}(\tfrac12,\tfrac12,-1)$, for which
$\mu_8^{\rm theirs} = (2/\sqrt3)\,\mu_8^{\rm ours} = 1.1547\,\mu_8^{\rm ours}$;
Buballa~\cite{Buballa2005} and Steiner, Reddy and Prakash~\cite{Steiner2002}
use the full $\lambda_8$, for which $\mu_8^{\rm ours} =
\sqrt3\,\mu_8^{\rm theirs}$. As a worked consequence, the CFL result
$\mu_8 = -\tfrac1{2\sqrt3} m_s^2/\mu$ of Ref.~\cite{Steiner2002} reads
$\mu_8 = -\tfrac12 m_s^2/\mu$ here: at $m_s = 300$, $\mu = 450$ MeV that is
$-57.7$ MeV in their convention and $-100.0$ MeV in ours.

\paragraph{Mode potentials.} With all five conserved-charge potentials,
\begin{equation}
  \mu_{(f,a)} = \frac{\muB}{3} + q_f \muC + s_f \muS
              + (T_3)_a\,\mu_3 + (T_8)_a\,\mu_8 ,
  \qquad \mu_e = -\muC ,
  \label{eq:modemu}
\end{equation}
the last of which is beta equilibrium in the sign convention
$\muC = \mu_p - \mu_n$ used across this repository.

\paragraph{Two naming traps in the sources.} The $G_D$ of
Ref.~\cite{Steiner2002} is the 't~Hooft coupling, not the diquark one (theirs
is $G_{\rm DIQ}$); and $\Delta_2$ in the R\"uster/Buballa index convention
means the \emph{ud} gap, not the second largest.

\section{The Lagrangian and what is kept}
\label{sec:lagrangian}

\begin{align}
  \mathcal{L} =\ & \bar q\,(i\slashed\partial - \hat m)\,q
   + \GS \sum_{a=0}^{8}\Big[ (\bar q\,\tau_a q)^2
                           + (\bar q\, i\gamma_5 \tau_a q)^2 \Big]
   \nonumber\\
   & - K \Big\{ \det\nolimits_f \big[\bar q (1+\gamma_5) q\big]
              + \det\nolimits_f \big[\bar q (1-\gamma_5) q\big] \Big\}
   \nonumber\\
   & + \GD \sum_{\eta}
        \big(\bar q\, i\gamma_5\, \epsilon_\eta \lambda^A_\eta\, C\, \bar q^{\,t}\big)
        \big(q^{t}\, C\, i\gamma_5\, \epsilon_\eta \lambda^A_\eta\, q\big)
   - \GV\,(\bar q \gamma^\mu q)^2 .
  \label{eq:lagrangian}
\end{align}
The scalar, vector and diquark mean fields all arise from \emph{one}
Hubbard--Stratonovich step: none is more fundamental than the others, and the
gap equation $\partial\Omega/\partial\Delta_\eta = 0$ and the mass equation are
the same kind of statement. Two places where the analogy nevertheless breaks
matter for the implementation. First, the condensation-energy normalisations
differ by a factor two --- the scalar cost is
$\sum_f (M_f-m_f)^2/(8\GS)$ and the pairing cost
$\sum_\eta \Delta_\eta^2/(4\GD)$ --- so $\eta_D = \GD/\GS = 1$ does
\emph{not} mean ``equally strong channels''. Second, $\sigma$ enters the
quasiparticle diagonal while $\Delta$ is strictly off-diagonal: the gap matrix
has identically zero diagonal, mixes particles with holes, and therefore needs
the doubled basis of Sec.~\ref{sec:bdg}. That is also why the gap kernel
carries branch signs where the mass equation does not.

\paragraph{Mean-field masses.} Including the determinant cross-terms, which are
the coefficients most often dropped or mis-signed,
\begin{equation}
  M_u = m_u - 4\GS \phi_u + 2K\,\phi_d \phi_s , \qquad
  M_d = m_d - 4\GS \phi_d + 2K\,\phi_u \phi_s , \qquad
  M_s = m_s - 4\GS \phi_s + 2K\,\phi_u \phi_d ,
  \label{eq:massgap}
\end{equation}
with $\phi_f = \langle \bar q_f q_f\rangle$, negative in the broken phase.

\paragraph{What is omitted.} The 't~Hooft term expanded in the presence of
diquark condensates generates a cross-term $\propto \sum_\alpha \sigma_\alpha
|\Delta_\alpha|^2$. Most treatments omit it; Baym \textit{et
al.}~\cite{Baym2018} include it with a coupling $K' \simeq K$. It is
\emph{not} implemented here, and consequently $\eta_D$ must be read as an
effective coupling that has absorbed it. Nothing here invents a coefficient by
analogy with the mass cross-terms.

\section{The cut medium integrals}
\label{sec:medium}

One mode of mass $M$ at effective potential $\mu^{*}$ and temperature $T$,
integrated to a cutoff $\Lam$, with $E = \sqrt{k^2+M^2}$ and the Fermi
functions $f^{\pm} = [1+e^{(E \mp \mu^{*})/T}]^{-1}$:
\begin{align}
  n     &= \frac{g}{2\pi^2}\int_0^{\Lam}\!\! dk\, k^2 \,(f^{+}-f^{-}) ,
  &
  \rho_s &= \frac{g}{2\pi^2}\int_0^{\Lam}\!\! dk\, k^2 \frac{M}{E}\,(f^{+}+f^{-}) ,
  \label{eq:medium1}\\
  \varepsilon &= \frac{g}{2\pi^2}\int_0^{\Lam}\!\! dk\, k^2\, E\,(f^{+}+f^{-}) ,
  &
  s &= \frac{g}{2\pi^2}\int_0^{\Lam}\!\! dk\, k^2 \sum_{\pm}
      \left[ \frac{x_{\pm}}{T} f^{\pm} + \ln\!\big(1+e^{-x_{\pm}/T}\big)\right] ,
  \label{eq:medium2}
\end{align}
with $x_{\pm} = E \mp \mu^{*}$, and the pressure
\begin{equation}
  P_{\log} = \frac{g}{2\pi^2}\int_0^{\Lam}\!\! dk\, k^2\, T
    \Big[ \ln\!\big(1+e^{-(E-\mu^{*})/T}\big)
        + \ln\!\big(1+e^{-(E+\mu^{*})/T}\big) \Big] .
  \label{eq:Plog}
\end{equation}
Antiparticles \emph{subtract} in $n$ and \emph{add} in $\rho_s$,
$\varepsilon$ and $P$. At $T = 0$ the occupations become step functions and
$T\ln(1+e^{-x/T}) \to \max(-x,0)$; $s$ vanishes identically.

\paragraph{The surface term.} The second standard pressure form,
\begin{equation}
  P_{k^4} = \frac{g}{6\pi^2}\int_0^{\Lam}\!\! dk\,
            \frac{k^4}{E}\,(f^{+}+f^{-}) ,
\end{equation}
is obtained from Eq.~\eqref{eq:Plog} by parts, and the boundary term does
\emph{not} vanish when the integral is cut:
\begin{equation}
  P_{\log} - P_{k^4} = \frac{g}{6\pi^2}\,\Lam^3\, T
    \Big[ \ln\!\big(1+e^{-(E_\Lam-\mu^{*})/T}\big)
        + \ln\!\big(1+e^{-(E_\Lam+\mu^{*})/T}\big) \Big] ,
  \qquad E_\Lam = \sqrt{\Lam^2+M^2} .
  \label{eq:surface}
\end{equation}
It is not small: $0.1\%$ of $P$ at $(M,\mu^{*},T) = (100,500,20)$ MeV,
$10.5\%$ at $(40,590,30)$, $36.3\%$ at $(140,700,5)$ and $39.9\%$ at
$(140,700,50)$. \emph{Every assembly in this implementation uses
$P_{\log}$.} At $T=0$ with $k_F < \Lam$ the two agree, which is exactly why
the error would hide until a table were built at finite temperature.

\paragraph{Quadrature.} Gauss--Legendre on one panel over $[0,\Lam]$ cannot
resolve the Fermi step. Breakpoints are placed at each Fermi momentum
$k_{F,j} = \sqrt{\mu^{*2}_j - M_{f(j)}^2}$ and, at $T>0$, at
$k_{F,j} \pm 25T$, and each panel is integrated separately; \emph{the cutoff is
imposed as the panel upper limit before the panels are built}, because
filtering breakpoints afterwards can delete the Fermi-surface break and revert
silently to single-panel accuracy.

\section{The Dirac sea, the condensates and the vacuum}
\label{sec:vacuum}

The vacuum integrals of one flavour, in closed form with $g_{\rm sea} = 6$ and
$R = \sqrt{\Lam^2+M^2}$:
\begin{align}
  \rho_s^{\rm vac}(M) &= \frac{g_{\rm sea}}{2\pi^2}\,\frac{M}{2}
      \Big[ \Lam R - M^2 \operatorname{arcsinh}(\Lam/M) \Big] ,
  \label{eq:seascalar}\\
  \varepsilon_{\rm sea}(M) &= \frac{g_{\rm sea}}{2\pi^2}\,\frac18
      \Big[ \Lam (2\Lam^2+M^2) R - M^4 \operatorname{arcsinh}(\Lam/M) \Big] ,
  \label{eq:seaeps}
\end{align}
and $\partial \varepsilon_{\rm sea}/\partial M = \rho_s^{\rm vac}$ identically.
The condensate of one flavour is the sea's scalar density with the medium's own
occupation subtracted, summed over the three colours of that flavour:
\begin{equation}
  \phi_f = -\rho_s^{\rm vac}(M_f) + \sum_{a} \rho_{s,(f,a)}
           + \delta\rho_{s,f} ,
  \label{eq:phi}
\end{equation}
where $\delta\rho_{s,f}$ is the pairing correction of
Sec.~\ref{sec:pairing}. The condensate cost entering both $\Omega$ and
$\varepsilon$ is
\begin{equation}
  C = 2\GS \sum_f \phi_f^2 - 4K\,\phi_u \phi_d \phi_s .
  \label{eq:C}
\end{equation}

\paragraph{The vacuum solution.} At $\mu = T = 0$ Eqs.~\eqref{eq:massgap} and
\eqref{eq:phi} close on themselves and are solved by a \emph{damped fixed
point} on the masses,
\begin{equation}
  M \leftarrow M + \lambda\,\big(M_{\rm new}[\phi(M)] - M\big), \qquad
  \lambda = 0.3 ,
\end{equation}
and not by a root finder on the condensates, which diverges and returns masses
that increase with density. The vacuum constant is
$\Omega_{\rm vac} = \varepsilon_{\rm vac}
= -\sum_f \varepsilon_{\rm sea}(M_f^{\rm vac}) + C^{\rm vac}$, and the equality
of the two is what makes the Euler relation survive the subtraction.

\paragraph{Vacuum diagnostics.} The pion decay constant from the quark loop,
\begin{equation}
  f_\pi^2 = 2 M^2 I_2 , \qquad
  I_2 = \frac{N_c}{4\pi^2}\Big[\operatorname{arcsinh}(\Lam/M)
        - \Lam/\sqrt{\Lam^2+M^2}\Big] ,
  \label{eq:fpi}
\end{equation}
and the effective bag constant, $\Omega$ at fixed masses evaluated at the
current masses minus at the broken-phase ones,
\begin{equation}
  B_{\rm eff} = \Omega\big[M_f = m_f\big] - \Omega\big[M_f = M_f^{\rm vac}\big] .
  \label{eq:Beff}
\end{equation}
$B_{\rm eff}$ is a \emph{derived} quantity here, not an input the way a bag
constant is in a bag model.

\section{The pairing sector}
\label{sec:pairing}

The machinery of this section is \texttt{eos/general/pairing.py}, shared with
the colour-dielectric model, because the pairing sector of the two is the same
sector.

\subsection{The gap matrix}
\label{sec:gapmatrix}

\begin{equation}
  G_{(fa),(gb)} = \sum_\eta \Delta_\eta\, \epsilon^{ab\eta}\,\epsilon_{fg\eta}
                = \sum_\eta \Delta_\eta\, (B_\eta)_{(fa),(gb)} ,
  \label{eq:G}
\end{equation}
so $\Delta_1$ pairs $d$ with $s$, $\Delta_2$ pairs $u$ with $s$, and
$\Delta_3$ pairs $u$ with $d$ --- the 2SC gap. The matrix is symmetric with
identically zero diagonal. Its eigenvalue multiplicities are a \emph{derived}
property of the pattern and are never assigned by hand; at $\Delta_0 = 60$ MeV,

\begin{center}
\begin{tabular}{lll}
\toprule
pattern & $\Delta$ & spectrum of $G$ [MeV] \\
\midrule
unpaired & $(0,0,0)$ & $0\,(\times 9)$ \\
2SC & $(0,0,\Delta_0)$ & $-60\,(\times2)$, $0\,(\times5)$, $+60\,(\times2)$ \\
CFL & $(\Delta_0,\Delta_0,\Delta_0)$ & $-60\,(\times5)$, $+60\,(\times3)$, $+120$ \\
uSC & $(0,\Delta_0,\Delta_0)$ & $\pm84.85$, $\pm60\,(\times2)$, $0\,(\times3)$ \\
dSC & $(\Delta_0,0,\Delta_0)$ & as uSC \\
\bottomrule
\end{tabular}
\end{center}
With independent gaps the $\pm\sqrt2\,\Delta_0$ eigenvalue generalises to
$\pm\sqrt{\Delta_2^2+\Delta_3^2}$.

\subsection{The Bogoliubov--de~Gennes problem}
\label{sec:bdg}

At each momentum $k$, for particles ($r=+$) and antiparticles ($r=-$)
separately, with $\xi^r_j = E_{f(j)} - r\,\mu^{*}_j$,
\begin{equation}
  H^r(k) = \begin{pmatrix}
     \mathrm{diag}(\xi^r) & G \\ G & -\mathrm{diag}(\xi^r)
  \end{pmatrix} \in \mathbb{R}^{18\times18} ,
  \label{eq:bdg}
\end{equation}
whose spectrum comes in $\pm$ pairs. The nine quasiparticle energies are the
\emph{non-negative half of the signed spectrum},
$E_a = \mathrm{sort}(\mathrm{eig}\,H)[9{:}]$ --- not the nine largest in
modulus. The two prescriptions agree in value but only the first is smooth
through a gapless window, where a branch crosses zero and its partner crosses
back, and that smoothness is what makes the Hellmann--Feynman derivatives below
carry the correct branch sign with no sign bookkeeping.

At unequal masses $[G, M] \neq 0$ --- $\|[G,M]\|_F = 7.4\times10^4$ at
$M = (40,45,480)$ MeV against exactly zero at equal masses --- so there is no
closed-form dispersion for a general pattern and Eq.~\eqref{eq:bdg} is
diagonalised numerically. The 2SC pattern is the exception:
\begin{equation}
  E^{\pm} = \sqrt{(\bar E - \bar\mu)^2 + \Delta^2} \pm
            \Big[\tfrac12 (E_d - E_u) - \delta\mu\Big] ,
  \qquad \bar E = \tfrac12 (E_u + E_d),\ \
  \bar\mu = \tfrac12(\mu^{*}_{ur}+\mu^{*}_{dg}),\ \
  \delta\mu = \tfrac12(\mu^{*}_{dg}-\mu^{*}_{ur}) ,
  \label{eq:twosc}
\end{equation}
which reproduces the numerical spectrum to $10^{-11}$ MeV over random
configurations and serves both as a fast path and as the unit test of the
general one. $E^{-}$ may be \emph{negative}: that is the gapless window, the
BCS blocking region, not an error.

\subsection{The pairing potential}
\label{sec:dOmega}

Written as a \emph{correction}, a difference from the unpaired spectrum:
\begin{equation}
  \delta\Omega_{\rm pair} = -\frac{1}{2\pi^2}\sum_{r=\pm}
    \int_0^{\Lam}\!\! dk\, k^2 \sum_{a=1}^{9}
    \Big[ \varphi(E^r_a) - \varphi(|\xi^r_a|) \Big] ,
  \qquad
  \varphi(x) = x + 2T \ln\!\big(1+e^{-x/T}\big) ,
  \label{eq:dOmega}
\end{equation}
with $\varphi(x) = x$ at $T = 0$. This vanishes \emph{identically} at
$\Delta = 0$, which is what makes the unpaired phase a clean limit of the same
code, and in the clean weak-coupling limit it obeys the BCS logarithm,
$-\delta\Omega_{\rm pair} \to (2/\pi^2)\,\mu^2\Delta^2[\ln(2\Lam/\Delta)-\tfrac12]$.
Both signs of $r$ are summed: the antiparticle branches contribute $8.8\%$ of
the particle piece at $\Lam = 600$ MeV and $17.1\%$ at $1000$ MeV.

The $|\xi_j|$ subtraction kinks at each of the nine $k_{F,j}$, so the pairing
quadrature is split there (and at $k_{F,j}\pm 25T$) exactly as in
Sec.~\ref{sec:medium}: splitting at $100$ nodes per panel reaches a relative
error of $3\times10^{-14}$ where a single panel of $800$ nodes reaches
$2\times10^{-7}$.

\subsection{The gap equations, and everything else the same pass returns}
\label{sec:kernels}

Since $\varphi'(x) = \tanh(x/2T)$, differentiating Eq.~\eqref{eq:dOmega} by
Hellmann--Feynman on Eq.~\eqref{eq:bdg} gives, with eigenvectors $|V_a\rangle$
in the doubled basis and $P_j$ the projector on mode $j$,
\begin{align}
  \frac{\Delta_\eta}{2\GD}
  - \frac{1}{2\pi^2}\sum_{r}\int\!\! dk\, k^2 \sum_a
    \Big\langle V^r_a \Big| \begin{pmatrix} 0 & B_\eta \\ B_\eta & 0\end{pmatrix}
    \Big| V^r_a \Big\rangle \tanh\!\frac{E^r_a}{2T} &= 0 ,
  \label{eq:gapeq}\\
  \delta n_j = \frac{1}{2\pi^2}\sum_{r}\int\!\! dk\, k^2
    \left[ \sum_a \tanh\!\frac{E^r_a}{2T}
      \frac{\partial E^r_a}{\partial \mu_j}
      + r \tanh\!\frac{\xi^r_j}{2T}\right] ,
  &\qquad
  \frac{\partial H^r}{\partial \mu_j} = -r\,(P_j \oplus -P_j) ,
  \label{eq:dn}\\
  \delta\rho_{s,f} = -\frac{1}{2\pi^2}\sum_{r}\int\!\! dk\, k^2
    \left[\sum_a \tanh\!\frac{E^r_a}{2T}
      \frac{\partial E^r_a}{\partial M_f}
      - \sum_{j \in f} \frac{M_f}{E_j}\tanh\!\frac{\xi^r_j}{2T}\right] ,
  &\qquad
  \frac{\partial H^r}{\partial M_f}
    = \mathrm{diag}(d) \oplus -\mathrm{diag}(d),\ \
    d_j = \frac{M_f}{E_j}\delta_{f(j),f} ,
  \label{eq:drhos}\\
  \delta s = \frac{1}{2\pi^2}\sum_{r}\int\!\! dk\, k^2 \sum_a
    \Big[ \psi(E^r_a) - \psi(|\xi^r_a|)\Big] ,
  &\qquad
  \psi(x) = 2\ln\!\big(1+e^{-x/T}\big)
            + \frac{x}{T}\Big[1 - \tanh\frac{x}{2T}\Big] .
  \label{eq:ds}
\end{align}
All four come from \emph{one} quadrature pass and one batched diagonalisation.

Three points about Eqs.~\eqref{eq:gapeq}--\eqref{eq:ds} are the ones that
return a plausible wrong answer if got wrong.

\paragraph{The kernel is not $\Delta/|E|$.} That form --- obtained by
differentiating $|E|$ as though every branch were positive --- is wrong by a
factor $12.0$ at $\Delta = 40$ MeV, $1.7$ at $60$ and $1.3$ at $80$ for
$\mu_u = 400$, $\mu_d = 500$ MeV, and it makes the gap \emph{grow} with the
mismatch, which is the opposite of the physics. It agrees with
Eq.~\eqref{eq:gapeq} only where every branch is positive.

\paragraph{Paired densities and entropy are not the unpaired integrals.} At
$\muB = 1400$ MeV, $T = 20$ MeV, $\Delta_3 = 80$ MeV the unpaired density
formula is wrong by $-21.1\%$ on the paired $u$ modes and $+11.6\%$ on the
paired $d$ modes. The entropy is worse: in a fully gapped phase the ratio of
paired to unpaired entropy is $2\times10^{-4}$ at $T = 5$ MeV.

\paragraph{The gap equation has three roots.} With a mismatch,
$R(\Delta) = \Delta/2\GD - \mathrm{kernel}$ vanishes at $\Delta = 0$, at a
barrier maximum, and at the physical BCS root. A fixed bracket returns
whichever it happens to contain. At $\mu^{*} = 450$ MeV, $\eta_D = 0.75$ the
roots are $92.71$ MeV at zero mismatch and $(32.4,\,92.71)$, $(52.5,\,92.71)$,
$(59.9,\,92.71)$ at $\delta\mu = 50, 60, 65$ MeV; the free-energy crossover
sits at $\delta\mu_c = 63.59$ MeV, a fraction $0.970$ of the weak-coupling
Clogston--Chandrasekhar value $\Delta_0/\sqrt2$, the $3\%$ deficit being the
finite cutoff.

\section{The vector sector}
\label{sec:vector}

With $\nq = \sum_j n_j$ the total quark density, the vector interaction energy
and self-energy are
\begin{equation}
  W(\nq) = \GV(\nq)\,\nq^2 , \qquad
  \Sigma_V = \frac{dW}{d\nq} = \big(2 - \alpha_{\rm eff}\big)\,\GV(\nq)\,\nq ,
  \qquad
  \alpha_{\rm eff} = -\frac{d\ln \GV}{d\ln \nq} ,
  \qquad
  \mu^{*}_j = \mu_j - \Sigma_V .
  \label{eq:vector}
\end{equation}
Three forms are implemented: constant, $\GV = \eta_V \GS$; a power law
$\GV = \GV^0 (n_{\rm ref}/\nq)^{\alpha}$; and the gluon-exchange form
\begin{equation}
  \GV(\nq) = \frac{\GV^0}{1 + 8 k_F^2/(9 M_g^2)} , \qquad
  k_F = \left(\tfrac12 \pi^2 \nq\right)^{1/3} ,
  \label{eq:gluonGV}
\end{equation}
for which $\alpha_{\rm eff} = \tfrac23\,u/(1+u)$ with $u = 8k_F^2/9M_g^2$.

\paragraph{Why the density dependence exists.} Writing
$\varepsilon = \sum_i C_i n^{p_i}$, each term contributes
$P_i = C_i (p_i-1) n^{p_i}$, so
\begin{equation}
  c_s^2 = \frac{\sum_i C_i p_i (p_i-1) n^{p_i-1}}
               {\sum_i C_i p_i n^{p_i-1}} ,
  \qquad\Longrightarrow\qquad
  c_s^2(n\to\infty) = \max\!\big(1-\alpha,\ \tfrac13\big)
\end{equation}
for $\GV \sim n^{-\alpha}$, since the vector term then has $p_V = 2-\alpha$
against the free-quark $4/3$. Constant $\GV$ gives $\alpha = 0$ and
$c_s^2 \to 1$, Zel'dovich behaviour. At $\alpha = 2/3$ the vector term is
conformal \emph{identically, not asymptotically}: $P_{\rm int}/\varepsilon_{\rm
int} = 1-\alpha = 1/3$ at every density. Eq.~\eqref{eq:gluonGV} reaches
$\alpha_{\rm eff} = 2/3$ as a consequence of its own structure, with no tuning
--- $\alpha_{\rm eff} = 0.062,\ 0.460,\ 0.608,\ 0.653$ at
$\nq = 10^6, 10^8, 10^9, 10^{10}$ MeV$^3$ --- which is why it is the
recommended variant. Pairing does not change the asymptotics:
$c_s^2 \to \tfrac13 + \tfrac29 \Delta^2/\mu^2$, which dies as $1/\mu^2$.

\paragraph{The rearrangement term is mandatory.} Once $\GV$ depends on the
density, $\Sigma_V$ is $dW/d\nq$ and not $2\GV \nq$; the ratio is $0.833$ at
$\alpha = 1/3$ and $0.667$ at $2/3$, and using the naive form shifts $P$ by
about $5\%$ and breaks $n = dP/d\mu$ at the first digit. (A coupling that
depends on a mean \emph{field} needs no such term, because it enters through
that field's own equation of motion. A density-dependent one does.)

\section{The totals}
\label{sec:totals}

With $D = \sum_\eta \Delta_\eta^2/(4\GD)$ the pairing cost and $C$ of
Eq.~\eqref{eq:C}, the matter sector alone --- no leptons, no photons ---
assembles as
\begin{align}
  \Omega &= -\sum_j P_{{\rm med},j} - \sum_f \varepsilon_{{\rm sea},f}
            + C - \big(\Sigma_V \nq - W\big)
            + \delta\Omega_{\rm pair} + D
            \ -\ \Omega_{\rm vac} ,
  \label{eq:Omega}\\
  \varepsilon &= \ \ \sum_j \varepsilon_{{\rm med},j}
            - \sum_f \varepsilon_{{\rm sea},f}
            + C + W + \varepsilon_{\rm pair}
            \ -\ \varepsilon_{\rm vac} ,
  \label{eq:eps}\\
  \varepsilon_{\rm pair} &= \delta\Omega_{\rm pair} + D
            + T\,\delta s + \sum_j \mu^{*}_j\, \delta n_j ,
  \qquad
  s = \sum_j s_{{\rm med},j} + \delta s ,
  \qquad
  P = -\Omega ,
  \label{eq:epspair}
\end{align}
and both carry the \emph{same} vacuum constant, so
$\Omega_{\rm vac} = \varepsilon_{\rm vac}$ exactly and the Euler relation
survives the subtraction. The conserved-charge sums are
\begin{equation}
  n_j = n_{{\rm med},j} + \delta n_j, \qquad
  \nq = \sum_j n_j , \qquad
  \nB = \frac{\nq}{3} , \qquad
  \nC = \sum_j q_{f(j)} n_j , \qquad
  \nS = \sum_j s_{f(j)} n_j ,
  \label{eq:charges}
\end{equation}
and the colour densities that neutrality sets to zero,
\begin{equation}
  n_3 = \sum_f \big(n_{(f,r)} - n_{(f,g)}\big) , \qquad
  n_8 = \sum_f \big(n_{(f,r)} + n_{(f,g)} - 2 n_{(f,b)}\big) ,
  \label{eq:colourdens}
\end{equation}
these being the generator densities up to the constants $\tfrac12$ and
$\tfrac13$; a row that must vanish does not care about its normalisation.

\paragraph{Euler.} Eqs.~\eqref{eq:Omega}--\eqref{eq:epspair} give
\begin{equation}
  \varepsilon + P = T s + \sum_j \mu_j n_j
  \label{eq:euler}
\end{equation}
with the \emph{physical} mode potentials, since
$\sum_j \mu^{*}_j n_j + \Sigma_V \nq = \sum_j \mu_j n_j$. This is audited at
every solved point. Three assembly bugs found during development each produced
a plausible equation of state and each was caught here: a sign error in
$\varepsilon$ ($\mathcal{O}(1)$), the pairing cost and $\delta\Omega_{\rm pair}$
dropped from both sums ($7.7\times10^{-3}$, small enough to pass for
quadrature), and $\delta s$ dropped, which fails only at $T>0$.

\paragraph{Leptons and thermal sectors.} Electrons at $\mu_e$, muons at
$\mu_\mu = \mu_e - \munue$, trapped neutrinos at $\munue$, photons, and the
untracked neutrino flavours as $\mu = 0$ gases, all from
\texttt{eos/general/thermodynamics\_leptons.py}. They are added to the totals
and are \emph{not} part of the phase: a phase does not own the leptons.

\section{The solve}
\label{sec:solve}

\subsection{The unknown vector}

\begin{equation}
  x = \Big(\underbrace{M_u, M_d, M_s}_{\text{always}},\ \
      \underbrace{\{\Delta_\eta\}_{\eta \in \text{pattern}}}_{\text{free gaps}},\ \
      \underbrace{\mu_3, \mu_8}_{\text{if the pattern pairs}},\ \
      \underbrace{\Sigma_V}_{\text{if } \GV \neq 0},\ \
      \underbrace{\muB, \muC}_{\text{always}},\ \
      \underbrace{\muS}_{\text{if } Y_S \text{ held}},\ \
      \underbrace{\munue}_{\text{if } Y_{L_e} \text{ held}} \Big) .
\end{equation}

\subsection{The rows, in the order the solver assembles them}

\begin{align}
  &M_f - \big[m_f - 4\GS \phi_f + 2K \phi_g \phi_h\big] = 0 ,
   && f = u,d,s \quad \text{(3 rows)} \\
  &\frac{\Delta_\eta}{2\GD} - \mathrm{kernel}_\eta = 0 ,
   && \text{one per free gap, Eq.~\eqref{eq:gapeq}} \\
  &n_3 = 0 , \qquad n_8 = 0 ,
   && \text{if the pattern pairs, Eq.~\eqref{eq:colourdens}} \\
  &\Sigma_V - \frac{dW}{d\nq} = 0 ,
   && \text{if } \GV \neq 0 \\
  &\nB - \nB^{\rm target} = 0 ,
   && \text{always} \\
  &\nC - n_e - n_\mu = 0 \ \ \text{or}\ \ \nC - Y_C \nB = 0 ,
   && \text{neutrality, or the held charge} \\
  &\nS - Y_S \nB = 0 ,
   && \text{if } Y_S \text{ held} \\
  &(n_e + n_{\nu_e}) - Y_{L_e}\nB = 0 ,
   && \text{if } Y_{L_e} \text{ held.}
\end{align}
Each row is divided by the scale of the quantity it balances --- $\muB$ for a
potential, $\nB$ for a density, $(\muB/3)^3/\pi^2$ for a colour row --- and the
state is accepted when the largest scaled component is below
$10^{-10}$. Without that scaling the norm would be dominated by the colour
rows, which are MeV$^3$, by twenty orders of magnitude.

\subsection{The modes}

\begin{center}
\begin{tabular}{llll}
\toprule
mode & variables & extra unknowns & rows replaced \\
\midrule
\texttt{beta\_eq\_neutrinoless} & $(\nB,T)$ & --- & neutrality; $\muS = 0$ \\
\texttt{beta\_eq\_neutrino\_trapped} & $(\nB, Y_{L_e}, T)$ & $\munue$ & neutrality $+$ $Y_{L_e}$ \\
\texttt{fixed\_YC} & $(\nB, Y_C, T)$ & --- & $\nC = Y_C \nB$; $\muS = 0$ \\
\texttt{fixed\_YC\_YS} & $(\nB, Y_C, Y_S, T)$ & $\muS$ & $\nC = Y_C \nB$, $\nS = Y_S \nB$ \\
\bottomrule
\end{tabular}
\end{center}
In a fixed-fraction mode with \texttt{leptons=True} the neutralizing leptons are
solved \emph{after} the matter, at the single potential that makes the total
system neutral: they feel no field the matter feels and nothing about them
feeds back. With \texttt{leptons=False} the result is charged matter, which is
what a mixed-phase construction needs per pure phase. Wherever a temperature is
accepted, the entropy per baryon may be given instead, as an outer
one-dimensional solve for $T$.

\subsection{The pattern is not a mode}
\label{sec:enumeration}

Which condensates survive is decided by free energy, not declared. Every
enumerated candidate --- unpaired, 2SC, CFL, and one asymmetric free seed that
can land on uSC, dSC or an unequal-gap state --- is solved to
self-consistency, and the converged one with the lowest $f = \varepsilon - Ts$
is returned; at fixed $\muB$ (which is what the phase adapter does) the
criterion is the largest $P$ instead, and the two agree. Every point reports
the winner, the three gaps, $\mu_3$, $\mu_8$ and whether the state is gapless.

Two seeding facts. CFL is electrically neutral \emph{without} electrons, so its
seed puts $\muC$ at zero; seeded with an electron-bearing potential the solve
reaches a spurious point with an $11\%$ flavour-density spread. And in an
unpaired region $\mu_8$ is unconstrained --- $n_8$ vanishes identically at
$\mu_8 = 0$ --- so it is pinned there rather than solved for.

\paragraph{The residual handed to the root finder is dimensionless.} The rows
above span twenty orders of magnitude in their raw units --- mass rows in MeV
against gap and colour rows in MeV$^3$ --- and a root finder terminates on its
own view of the residual, not on the scaled one the gate is stated in. So each
row is divided by its scale \emph{before} the vector is handed over, with a
matching tolerance. Without that, every row but the largest is driven only as
far as the largest one needs, and a paired solve reports $10^{-8}$ where it can
reach $10^{-16}$; note in particular that a gap row is a \emph{density}, since
$\Delta_\eta/2\GD$ carries MeV$^3$, and judging it against a potential is four
orders of magnitude too strict.

\section{The phase-adapter surface}
\label{sec:adapter}

\texttt{eos/mixed} consumes this model through one function: given
$(\muB, \muC, \muS, T)$ it returns the phase block, having closed the model's
own internal system --- masses, gaps, $\Sigma_V$ \emph{and the two colour
potentials}. Colour neutrality is internal because $\mu_3$ and $\mu_8$ are not
conserved charges of the mixed system: no hadronic phase carries them and there
is nothing across the interface for them to equilibrate with. The slot carries
the \emph{physical} baryon potential, and the seed is not cacheable: the seed
chooses the root, so caching it would change physics rather than speed.

\section{What the implementation reproduces}
\label{sec:reproduces}

\begin{center}
\begin{tabular}{lll}
\toprule
quantity & computed & reference \\
\midrule
$M_u$ (vacuum) & $367.648$ MeV & $367.7$~\cite{Rehberg1996} \\
$M_s$ (vacuum) & $549.479$ MeV & $549.5$ \\
$(-\phi_u)^{1/3}$ & $241.946$ MeV & $241.9$ \\
$(-\phi_s)^{1/3}$ & $257.688$ MeV & $257.7$ \\
$f_\pi$ & $92.391$ MeV & $92.4$ \\
$B_{\rm eff}^{1/4}$ & $228.93$ MeV & $357.49$ MeV\,fm$^{-3}$ \\
2SC gap at $\mu^{*}=450$ MeV, $\eta_D = 0.75$ & $92.71$ MeV & --- \\
Clogston ratio $\delta\mu_c/(\Delta_0/\sqrt2)$ & $0.970$ & $1$ (weak coupling) \\
\midrule
\multicolumn{3}{l}{\emph{Neutral solved points at $\muB = 1500$ MeV, $T = 0$,
$\eta_D = 0.75$:}} \\
unpaired & $M = (9.84, 8.55, 265.59)$ MeV, $\muC = -34.20$,
           $\nB = 1.4319$ fm$^{-3}$, $P = 302.12$ MeV\,fm$^{-3}$ \\
2SC & $\Delta_3 = 95.50$, $\mu_3 = 0$, $\mu_8 = -2.46$, $M_s = 243.13$ MeV,
      $\muC = -62.27$, $\nB = 1.4887$ fm$^{-3}$, $P = 324.75$ MeV\,fm$^{-3}$ \\
\bottomrule
\end{tabular}
\end{center}

One deliberate discrepancy is recorded rather than reconciled. The
implementation specification reports $M_u, M_d = (11.96, 7.65)$ MeV for the
2SC point above and this implementation gives $(9.73, 8.90)$. The difference is
the sign of the hole amplitudes in the paired scalar density $\delta\rho_{s,f}$
of Eq.~\eqref{eq:drhos}, and the identity $\nB = dP/d\muB$ along the neutral
solution decides it: it holds to the finite-difference floor with the sign used
here and fails by $2.6\times10^{-4}$ with the other. Near chiral restoration
$M-m$ is the small difference of two large numbers, which is why a
percent-level change in one scalar density moves the light masses by twenty
percent and moves $M_s$, $\Delta_3$, $\mu_8$, $\muC$, $\nB$ and $P$ not at all.

\section{What is not implemented}
\label{sec:notimplemented}

Recorded in \texttt{docs/DEFERRED.md} rather than approximated:
the RG-consistent regularization at $\lambda = \Lam_{\rm UV}/\Lam > 1$ (the
parameter exists and any value but $1$ raises, since the counterterm
cancelling the medium's logarithmic divergence $-(2/\pi^2)\bar\mu^2\Delta^2
\ln\Lam_{\rm UV}$ is not written); the 't~Hooft--diquark cross-term $K'$ of
Sec.~\ref{sec:lagrangian}; the trapped muon lepton family as a conserved
charge; the composition and gap freezes of \texttt{eos\_response}; and the
dilaton/colour-dielectric graft, which the specification marks as unverified
for transition order, pairing coexistence and finite temperature.

\begin{thebibliography}{9}
\bibitem{Rehberg1996} P.~Rehberg, S.~P.~Klevansky, J.~H\"ufner,
  Phys.\ Rev.\ C \textbf{53}, 410 (1996), arXiv:hep-ph/9506436.
\bibitem{Buballa2005} M.~Buballa, Phys.\ Rept.\ \textbf{407}, 205 (2005),
  arXiv:hep-ph/0402234.
\bibitem{Ruester2005} S.~B.~R\"uster, V.~Werth, M.~Buballa, I.~A.~Shovkovy,
  D.~H.~Rischke, Phys.\ Rev.\ D \textbf{72}, 034004 (2005),
  arXiv:hep-ph/0503184.
\bibitem{Steiner2002} A.~W.~Steiner, S.~Reddy, M.~Prakash,
  Phys.\ Rev.\ D \textbf{66}, 094007 (2002), arXiv:hep-ph/0205201.
\bibitem{Alford2008} M.~G.~Alford, A.~Schmitt, K.~Rajagopal, T.~Sch\"afer,
  Rev.\ Mod.\ Phys.\ \textbf{80}, 1455 (2008), arXiv:0709.4635.
\bibitem{Baym2018} G.~Baym, T.~Hatsuda, T.~Kojo, P.~D.~Powell, Y.~Song,
  T.~Takatsuka, Rept.\ Prog.\ Phys.\ \textbf{81}, 056902 (2018),
  arXiv:1707.04966.
\bibitem{Kunkel2026} S.~Kunkel, I.~A.~Rather \textit{et al.},
  arXiv:2607.11537.
\bibitem{Pagliara2008} G.~Pagliara, J.~Schaffner-Bielich,
  Phys.\ Rev.\ D \textbf{77}, 063004 (2008), arXiv:0711.1119.
\bibitem{Gholami2024} H.~Gholami, M.~Hofmann, M.~Buballa,
  arXiv:2408.06704.
\end{thebibliography}

\end{document}
